\section{A genus-one curve with large index and independent labels}
\label{sec:twist}

The fields in this section are auxiliary to the block coefficient system.
Fix a prime $p$ and put $\kappa=\overline{\mathbb F}_p$.  For $n=p^s$ we
will construct a genus-one curve $E/F_0$ with a finite \'{e}tale divisor
$\mathcal B$ of degree $n$. Every divisor on $E$ will have degree
divisible by $n$. We will also choose a constant elliptic curve
$A_0/\kappa$ and a cyclic extension $F'/F_0$ with
$F'=\kappa(A_0^n)$, over which the marked points identify with the
generic projection points $t_i$ of $A_0^n$. For a nonzero invariant
differential $\eta$ on $A_0$, the pullbacks $t_i^*\eta$ will form a
basis of $\Omega_{F'/\kappa}$. These are absolute differentials over
$\kappa$, not relative differentials over $F'$: the coordinates of
the marked points are themselves parameters in $F'$. We will bound
how much the divisor index and the rank of these forms can decrease
after finite extensions whose degree has small $p$-part.

For context, arbitrary-index genus-one curves over infinite fields
finitely generated over their prime fields are constructed by Clark and
Lacy \cite{ClarkLacy2019}. The fields here are instead finitely
generated over $\kappa$, and our curve must carry the independent
marked differential labels as well as the prescribed index. We prove
both properties for one translated cyclic twist.

\subsection{An elliptic curve with arbitrarily deep torsion}

An elliptic curve in characteristic $p$ is ordinary precisely when it
has nonzero geometric $p$-torsion.  We construct this property directly.

\begin{lemma}\label{lem:ordinary-elliptic}
For every prime $p$ there is an elliptic curve $A_0/\kappa$ with a point
of order $p$.  On this same curve, for every $s\geq 1$ there is a point
$P\in A_0(\kappa)$ of exact order $p^s$.
\end{lemma}

\begin{proof}
For $p=2$, take the Weierstrass cubic
\[
                     y^2+xy=x^3+1.
\]
The affine partial derivatives could vanish simultaneously only at
$(x,y)=(0,0)$, which is not on the curve.  In its projective equation
$Y^2Z+XYZ=X^3+Z^3$, the origin $[0:1:0]$ is smooth because the partial
derivative with respect to $Z$ is nonzero there.  Its points over
$\mathbb F_2$ are
\[
              [0:1:0],\quad (0,1),\quad (1,0),\quad (1,1).
\]
Thus its finite group of rational points contains an element of order
$2$, by Sylow theory.

Suppose that $p$ is odd and put $r=(p-1)/2$.  Consider
\[
 f_\lambda(x)=x(x-1)(x-\lambda),\qquad
 H(\lambda)=[x^{p-1}]f_\lambda(x)^r.
\]
Here $[x^a]$ means the coefficient of $x^a$.  The coefficient of
$\lambda^r$ in $H$ is $(-1)^r$: in
$x^r(x-1)^r(x-\lambda)^r$, take $(-\lambda)^r$ from the last factor
and $x^r$ from the middle factor.  In particular $H$ is not the zero
polynomial.  Choose $\lambda\in\kappa\setminus\{0,1\}$ with
$H(\lambda)\ne0$.  The cubic
\[
                         A_0:\quad y^2=f_\lambda(x)
\]
is smooth, since $f_\lambda$ has three distinct roots and the point at
infinity is smooth.

Choose $q=p^\ell$ with $\lambda\in\mathbb F_q$.  For each
$x\in\mathbb F_q$, the number of solutions for $y$ is
$1+\chi_q(f_\lambda(x))$, where $\chi_q$ is the quadratic character
with values in $\{0,1,-1\}$.  Its reduction in $\mathbb F_q$ is
$f_\lambda(x)^{(q-1)/2}$.  Reducing the point count modulo $p$
therefore gives
\begin{equation}\label{eq:elliptic-point-count}
 \#A_0(\mathbb F_q)
 \equiv 1+\sum_{x\in\mathbb F_q}f_\lambda(x)^{(q-1)/2}
 =1-H(\lambda)^{1+p+\cdots+p^{\ell-1}}.
\end{equation}
We verify the last equality, including its coefficient calculation.
The exponents occurring in $f_\lambda(x)^{(q-1)/2}$ range from
$(q-1)/2$ to $3(q-1)/2$.  Among them, only $q-1$ is a positive
multiple of $q-1$.  Since $\mathbb F_q^\times$ is cyclic, the sum of
$x^a$ over $\mathbb F_q$ is zero for the other positive exponents and
is $-1$ for $a=q-1$.
Furthermore,
\[
 f_\lambda(x)^{(q-1)/2}
       =\prod_{j=0}^{\ell-1}\bigl(f_\lambda(x)^r\bigr)^{p^j}.
\]
A contribution to the coefficient of $x^{q-1}$ would choose exponents
$a_j\in[r,3r]$ satisfying
\[
                   \sum_{j=0}^{\ell-1}p^j a_j=p^\ell-1.
\]
Reduction modulo $p$ forces $a_0\equiv p-1\pmod p$.  The interval
$[r,3r]$ has length $p-1$ and contains exactly one such integer,
namely $p-1$.  Subtract this integer and divide the displayed equality
by $p$.  Repetition forces $a_j=p-1$ for every $j$.  Thus the desired
coefficient is exactly
$\prod_jH(\lambda)^{p^j}$, proving
\eqref{eq:elliptic-point-count}.

The nonzero element
$c=H(\lambda)^{1+p+\cdots+p^{\ell-1}}$ lies in
$\mathbb F_p^\times$.  Replace $\mathbb F_q$ by
$\mathbb F_{q^d}$, where $d$ is a multiple of the multiplicative order
of $c$.  The corresponding expression is $c^d=1$, so
\eqref{eq:elliptic-point-count} shows that
$p\mid\#A_0(\mathbb F_{q^d})$.  Sylow theory again supplies a point
of order $p$.

Finally, multiplication by $p$ is a finite locally free map of degree
$p^2$ on an elliptic curve, by \Cref{foundation:elliptic}.  It is
surjective, hence surjective on points over the algebraically closed
field $\kappa$.  If $Q$ has order $p^a$ and $[p]Q'=Q$, then
$[p^{a+1}]Q'=0$ whereas $[p^a]Q'\ne0$.  Starting from a point of
order $p$ and repeatedly choosing such a preimage proves the assertion.
\end{proof}

\subsection{Cycle sums and rational maps to a constant elliptic curve}

The twist below uses $n=p^s$, a point $P\in A_0(\kappa)$ of exact
order $n$, and the field $F'=\kappa(A_0^n)$ with generic projection
points $t_i$. Its field automorphism and semilinear descent action
on $A_{0,F'}$ will be
\[
 \sigma(t_i)=t_{i+1}-P,\qquad \delta(x)=\sigma(x)+P
 \qquad(i\text{ modulo }n).
\]
The opposite translations make $\delta$ permute the marked points.
To test the degree $b$ of a divisor on the descended curve, we will
pull it back to $A_{0,F'}$ and take its geometric group-law sum $v$.
Invariance of the divisor under $\delta$ will give
\[
                         \sigma(v)+[b]P=v.
\]
The first lemma shows that $v$ is rational over $F'$, including when
the divisor has inseparable residue fields. Such a point is a rational
map $A_0^n\dashrightarrow A_0$. The second lemma writes it as a
constant plus one endomorphism in each coordinate. This will let us
compare the coordinate endomorphisms in the displayed descent relation
and determine which degrees $b$ are possible.

\begin{lemma}\label{lem:cycle-sum-rationality}
Let $A/\kappa$ be an elliptic curve and $U/\kappa$ a field extension.
For every zero-cycle $z$ on $A_U$, its group-law sum over an algebraic
closure of $U$, using the full geometric multiplicities, belongs to
$A(U)$.  This sum is additive for signed cycles.  Translation of a
cycle of degree $b$ by a point $Q\in A(U)$ adds $[b]Q$ to its sum.
\end{lemma}

\begin{proof}
For a point with inseparable residue field, its multiplicity in the
geometric cycle must be included before applying Galois descent to
the sum. We first show that multiplication by a power of $p$ moves
the point into the corresponding field of $p$th powers.
Write $K_A=\kappa(A)$.  A nonzero invariant differential $\eta$
generates $\Omega_{K_A/\kappa}$, and $[p]^*\eta=0$ by
\Cref{foundation:elliptic}.  Thus every rational function pulled back
by $[p]$ has differential zero.  The kernel of
$d:K_A\longrightarrow\Omega_{K_A/\kappa}$ is exactly $K_A^p$.
Indeed, it is a subfield containing $K_A^p$, it is a proper subfield
because $\Omega_{K_A/\kappa}$ has dimension one, and
$[K_A:K_A^p]=p$ by \Cref{foundation:fields}.  Consequently
\begin{equation}\label{eq:multiplication-p-powers}
                         [p]^*K_A\subseteq K_A^p.
\end{equation}

For every extension $V/\kappa$, this entails
\begin{equation}\label{eq:multiplication-point-fields}
                         [p]\bigl(A(V)\bigr)\subseteq A(V^p).
\end{equation}
Here $\kappa\subseteq V^p$ because $\kappa$ is perfect.  To justify
evaluation in \eqref{eq:multiplication-point-fields}, take
$R\in A(V)$ and a $\kappa$-affine open $W$ containing $[p]R$.
For $f\in\mathcal O_A(W)$, write $[p]^*f=g^p$ with
$g\in K_A$, using \eqref{eq:multiplication-p-powers}.  On
$[p]^{-1}W$ the rational function $g$ is regular: its $p$th power is
regular, and each local ring of this smooth curve is integrally
closed.  Hence $f([p]R)=g(R)^p\in V^p$ for every such $f$.
The point $[p]R$ therefore factors through $W(V^p)$, as claimed.
Iteration gives
$[p^u](A(V))\subseteq A(V^{p^u})$.

It suffices to treat a closed point $z$ with residue field $T/U$.
Let $T_{\mathrm{sep}}$ be the maximal separable subextension and write
$[T:T_{\mathrm{sep}}]=p^u$.  Then
$T^{p^u}\subseteq T_{\mathrm{sep}}$.  Over an algebraic closure
$\overline U$, the cycle of $z$ consists of the points $R_\tau$
indexed by the $U$-embeddings
$\tau:T_{\mathrm{sep}}\hookrightarrow\overline U$, each with
multiplicity $p^u$.  To see the multiplicity, first split the separable
extension; each remaining purely inseparable tensor factor is local
of dimension $p^u$ over $\overline U$, with residue field
$\overline U$.

Each embedding $\tau$ extends uniquely to $T$ inside $\overline U$.
By \eqref{eq:multiplication-point-fields},
\[
             [p^u]R_\tau\in A\bigl(\tau(T_{\mathrm{sep}})\bigr).
\]
The geometric group-law sum is therefore
$\sum_\tau[p^u]R_\tau\in A(U^{\mathrm{sep}})$, and Galois
permutes its summands.  Galois descent for points gives a point of
$A(U)$.  Additivity proves the assertion for arbitrary signed cycles.
Finally, in the geometric sum each copy of each point gains $Q$ under
translation, so the change is exactly $[b]Q$.
\end{proof}

\begin{lemma}\label{lem:constant-product-maps}
Let $A/\kappa$ be an elliptic curve.
\begin{enumerate}[label=\textup{(\roman*)}]
\item For every field extension $V/\kappa$, an origin-preserving
morphism $A_V\longrightarrow A_V$ is a group endomorphism defined
over $\kappa$.
\item Every rational map $A^n\dashrightarrow A$ has a unique expression
\begin{equation}\label{eq:constant-product-map}
 (x_1,\ldots,x_n)\longmapsto
                  c+\sum_{i=1}^n\beta_i(x_i),
 \qquad c\in A(\kappa),\quad
        \beta_i\in\operatorname{End}_\kappa(A).
\end{equation}
In particular, it extends to a morphism on the whole product.
\end{enumerate}
\end{lemma}

\begin{proof}
We first prove the homomorphism assertion in (i).  It can be checked
after extending $V$ to an algebraic closure.  If $f(0)=0$, form
\[
            g(x,y)=f(x+y)-f(x)-f(y).
\]
Choose a point $a\ne0$ of the target elliptic curve.  The projection
to the second factor of $g^{-1}(a)$ is closed, by properness, and
does not contain $0$, because $g(x,0)=0$.  Its complement is a
nonempty open set $W$.  For $y\in W$, the morphism
$x\mapsto g(x,y)$ misses $a$.  A nonconstant morphism of projective
integral curves is surjective, so this morphism is constant; its value
at $x=0$ is zero.  Thus $g$ is zero on $A\times W$, and hence
everywhere.  This proves that $f$ is a homomorphism.

We next descend the homomorphism to $\kappa$, first by controlling its
values on $A(\kappa)$ and then by recovering its rational functions
from those values. The zero homomorphism already descends, so suppose
that $f$ is nonconstant. Because $\kappa=\overline{\mathbb F}_p$, the
curve, its origin, and any chosen point of $A(\kappa)$ are defined over
a common finite field. The curve has a finite group of points over
that field, so every point of $A(\kappa)$ is torsion. Since $f$ is a homomorphism, each
such point has torsion image. For each positive integer $N$, the scheme
$A[N]$ is finite over
the algebraically closed field $\kappa$; all its points with values
in an extension field are therefore $\kappa$-points.  Consequently
\begin{equation}\label{eq:constant-endomorphism-values}
                           f(A(\kappa))\subseteq A(\kappa).
\end{equation}

For a rational function $h\in\kappa(A)$, write its pullback as
\[
 f^*h=\frac{a}{b},\qquad
 a=\sum_{j=1}^t c_j a_j,\quad
 b=\sum_{j=1}^t c_j b_j,
 \qquad a_j,b_j\in\kappa(A),
\]
where the finitely many $c_j\in V$ are linearly independent over
$\kappa$.  Such an expression follows by putting a rational function
in $V(A)$ over a common denominator and choosing a basis for its
finitely many scalar coefficients.  Fix $j_0$ with $b_{j_0}\ne0$.
For all but finitely many $R\in A(\kappa)$, all functions in this
expression can be evaluated, $b(R)b_{j_0}(R)\ne0$, and
$f^*h(R)\in\kappa$, by \eqref{eq:constant-endomorphism-values}.
Writing $\lambda_R=f^*h(R)$, scalar independence gives
\[
       a_j(R)=\lambda_R b_j(R)\quad\text{for every }j,
       \qquad
       f^*h(R)=\frac{a_{j_0}(R)}{b_{j_0}(R)}.
\]
There are infinitely many such $\kappa$-points.  A nonzero rational
function on a projective integral curve has only finitely many zeros,
so $f^*h=a_{j_0}/b_{j_0}\in\kappa(A)$.
Apply this to coordinate functions on a $\kappa$-affine open of the
target.  The morphism $f$ is thus a rational map defined over $\kappa$.
It extends over every point of the smooth source curve by the
valuative criterion for properness, and the extension agrees with $f$
after scalar extension.  This proves (i).

For (ii), view a rational map in the first variable over
$V=\kappa(A^{n-1})$, the function field of the remaining factors.
It extends to a morphism $A_V\to A_V$: the local rings at closed
points of the smooth source curve are DVRs, and properness extends
the map uniquely over each of them.  These extensions are morphisms
on neighborhoods and glue by separatedness.  Subtracting its value
at the origin leaves, by (i), an endomorphism
$\beta_1\in\operatorname{End}_\kappa(A)$.  The subtracted value is
an element of $A(V)$, hence a rational map from the remaining
$n-1$ factors.  Repeat to obtain \eqref{eq:constant-product-map}.
The right-hand side is a morphism everywhere.  Its value at the
origin determines $c$, and its restriction to each factor with the
other coordinates zero determines the corresponding $\beta_i$.
This also proves uniqueness.
\end{proof}

\subsection{The translated cyclic twist and its index}

For a regular projective integral curve over a field, its
\emph{index} is the positive generator of the subgroup of $\mathbb Z$
consisting of divisor degrees; equivalently, it is the greatest common
divisor of the degrees of its closed points.  Degrees throughout are
taken over the stated field.

\begin{proposition}\label{prop:twist-index}
Fix $s\geq1$, put $n=p^s$, and choose $A_0/\kappa$ and
$P\in A_0(\kappa)$ of exact order $n$ as in
\Cref{lem:ordinary-elliptic}.  There are finitely generated fields
$F_0\subset F'$ over $\kappa$, each of transcendence degree $n$,
and a smooth projective geometrically integral genus-one curve
$E/F_0$ with the following properties:
\begin{enumerate}[label=\textup{(\roman*)}]
\item $F'=\kappa(A_0^n)$, and $F'/F_0$ is cyclic Galois of degree
$n$.  With $t_i\in A_0(F')$ the generic projection points, a
generator satisfies
\begin{equation}\label{eq:shifted-cyclic-action}
                 \sigma(t_i)=t_{i+1}-P
                 \qquad(i\text{ modulo }n).
\end{equation}
\item Under an identification $E_{F'}\simeq A_{0,F'}$, the descent
action on points is $x\mapsto\sigma(x)+P$.  The two disjoint divisors
\begin{equation}\label{eq:twist-divisors}
       \mathcal B_{F'}=\sum_{i=1}^n[t_i],\qquad
       D_{F'}=\sum_{j=0}^{n-1}[jP]
\end{equation}
descend to finite \'{e}tale effective divisors $\mathcal B,D$ of
degree $n$ on $E$, and $D$ is ample.
\item The index of $E/F_0$ is exactly $n$.
\end{enumerate}
\end{proposition}

\begin{proof}
The cyclic permutation followed by translation on the smooth product
$A_0^n$ defines the automorphism of its function field in
\eqref{eq:shifted-cyclic-action}.  Its powers satisfy
\[
                         \sigma^a(t_i)=t_{i+a}-[a]P.
\]
Thus $\sigma^n=1$.  If $0<a<n$, the maps given by the independent
projections $t_i$ and $t_{i+a}$ cannot differ by a constant: vary one
factor while holding the other fixed.  Therefore $\sigma^a\ne1$.
Set $F_0=(F')^{\langle\sigma\rangle}$.  Finite fixed-field theory
gives $[F':F_0]=n$ with the indicated cyclic Galois group.  Since
$F'/\kappa$ is finitely generated of transcendence degree $n$, the
same is true of $F_0/\kappa$.  More explicitly, choose a
transcendence basis of $F_0/\kappa$; it is also a transcendence basis
of $F'/\kappa$, and $F'$ is finite over the rational function field
on this basis, as is its intermediate field $F_0$.

On $A_{0,F'}$ use the semilinear action
\[
                             \delta(x)=\sigma(x)+P.
\]
Since $\sigma$ fixes $P$ and $[n]P=0$, this action satisfies
$\delta^n=1$. The translation by $-P$ in the field action cancels
the translation by $P$ here, so $\delta$ sends $t_i$ to $t_{i+1}$.
It also sends $jP$ to $(j+1)P$. Thus the generic projections give the
invariant marked divisor $\mathcal B_{F'}$, and the orbit of the
origin gives the separate invariant divisor $D_{F'}$ of degree $n$.
All their points are distinct, and the two supports are
disjoint: a generic projection is nonconstant, so cannot equal a
constant point $jP$.

We use $D_{F'}$ to descend the curve projectively; once descended,
its degree will also give the upper bound $n$ for the index.
The invertible sheaf $\mathcal O(D_{F'})$ has a canonical
linearization, obtained by applying the action to rational functions
with prescribed poles in the invariant divisor.  It is ample: its
third power is the tensor product of the translates of
$\mathcal O(3[0])$ by the points $jP$, and $\mathcal O(3[0])$ is
the ample Weierstrass hyperplane bundle.  Thus projective Galois
descent, in the form of \Cref{foundation:descent}, constructs $E$
and descends both divisors.  Smoothness, geometric integrality,
finiteness, and the \'{e}tale property follow after the splitting
extension; ampleness descends as well.  Cohomology and base change
give $\dim_{F_0}H^1(E,\mathcal O_E)=1$, so $E$ has genus one.

It remains to prove that every divisor degree is divisible by $n$.
The two preparatory lemmas reduce this to comparing endomorphisms in
the cyclic descent relation.
Let $Z$ be an arbitrary divisor on $E$ of degree $b$, and pull it
back to $A_{0,F'}$.  Pullback preserves its degree.  By
\Cref{lem:cycle-sum-rationality}, its geometric group-law sum,
including inseparable multiplicities and signed coefficients, is a
point $v\in A_0(F')$.  The pulled-back divisor is invariant under
$\delta$.  Applying the translation formula in that lemma yields
\begin{equation}\label{eq:divisor-descent-sum}
                               \sigma(v)+[b]P=v.
\end{equation}
By \Cref{lem:constant-product-maps}, the point $v$, viewed as a
rational map from $A_0^n$ to $A_0$, has the form
\[
                  v=c+\sum_{i=1}^n\beta_i(t_i),\qquad
                  c\in A_0(\kappa),\quad
                  \beta_i\in\operatorname{End}_\kappa(A_0).
\]
Substitution in \eqref{eq:divisor-descent-sum} gives
\begin{equation}\label{eq:index-endomorphism-comparison}
 \sum_{i=1}^n(\beta_{i-1}-\beta_i)(t_i)
           +[b]P-\sum_{i=1}^n\beta_i(P)=0,
 \qquad \beta_0=\beta_n.
\end{equation}
This equality at the generic point is an equality of morphisms on
the product.  Evaluate it first at the origin and then on one factor
at a time.  Each $\beta_{i-1}-\beta_i$ must be zero, so all
$\beta_i$ equal one endomorphism $\beta$.  Their translation terms
sum to
\[
                         \sum_i\beta_i(P)=\beta([n]P)=0.
\]
Equation \eqref{eq:index-endomorphism-comparison} now gives
$[b]P=0$.  Since $P$ has exact order $n$, this proves $n\mid b$.
The divisor $D$ has degree $n$, so the index is exactly $n$.
\end{proof}

The marked divisor also has the differential property stated at the
start of the section. Fix a nonzero invariant differential $\eta$ on
$A_0$. The product formula for K\"ahler differentials, evaluated at the
generic point of $A_0^n$, gives
\begin{equation}\label{eq:label-differential-basis}
 \Omega_{F'/\kappa}=\bigoplus_{i=1}^n F'\,t_i^*\eta.
\end{equation}
Thus the same splitting field that supplies the marked points makes
their absolute differential forms independent.

\subsection{Consequences after finite extension}

\begin{corollary}\label{cor:index-base-change}
Let $h\geq1$ and $s\geq h$, and let $F/F_0$ be a finite extension
with $[F:F_0]_p\leq p^h$.  Every divisor on $E_F$ has degree over
$F$ divisible by $p^{s-h}$.  More generally, if $C/F$ is a regular
projective integral curve with a finite dominant morphism
$C\to E_F$, then every divisor and every invertible sheaf on $C$
has degree over $F$ divisible by $p^{s-h}$.
\end{corollary}

\begin{proof}
Put $N=[F:F_0]$.  For a closed point $z$ of $E_F$ mapping to $x$
of $E$, finite pushforward is
$z\mapsto[\kappa(z):\kappa(x)]x$.  Its degree over $F_0$ is
\[
 [\kappa(z):\kappa(x)]\,[\kappa(x):F_0]
       =[\kappa(z):F_0]=N[\kappa(z):F].
\]
Thus a divisor of degree $b$ over $F$ pushes forward to a divisor
of degree $Nb$ over $F_0$, including when $F/F_0$ or a point's
residue extension is inseparable.  By \Cref{prop:twist-index},
$p^s\mid Nb$.  Since $N_p\leq p^h$, it follows that
$p^{s-h}\mid b$.

For $C\to E_F$, the same residue-field tower calculation shows
that pushforward preserves the degree over $F$.  Apply the first
assertion to this pushforward.  Finally every invertible sheaf on a
regular integral curve has a nonzero rational section and is the
invertible sheaf of its divisor, so the line-bundle assertion follows.
\end{proof}

\begin{lemma}\label{lem:label-rank}
Let $F/F_0$ be a finite extension with $[F:F_0]_p\leq p^h$, and
choose a compositum $S=FF'$.  Then $S/F$ is finite Galois of degree
dividing $n$, and
\begin{equation}\label{eq:compositum-small-p-part}
                            [S:F']_p\leq p^h.
\end{equation}
Fix a nonzero invariant differential $\eta$ on $A_0$.  For every
subset $I\subseteq\{1,\ldots,n\}$, the forms $t_i^*\eta$,
$i\in I$, have rank at least $|I|-h$ in $\Omega_{S/\kappa}$.
The covectors
\begin{equation}\label{eq:label-covectors}
                    (1,-t_i^*\eta)\in S\oplus\Omega_{S/\kappa},
                    \qquad i\in I,
\end{equation}
also have rank at least $|I|-h$.
\end{lemma}

\begin{proof}
Since $F'/F_0$ is Galois, its compositum with any finite extension
$F/F_0$, including an inseparable one, is Galois over $F$, with
Galois group a subgroup of $\operatorname{Gal}(F'/F_0)$.
Consequently $[S:F]$ divides $n=p^s$.  The tower identity
\[
                    [S:F']\,[F':F_0]=[S:F]\,[F:F_0]
\]
then gives \eqref{eq:compositum-small-p-part}.

The forms $t_i^*\eta$ are the basis in
\eqref{eq:label-differential-basis}. Their loss of independence over $S$ is measured
by the kernel of the first map in the exact sequence
\[
 S\otimes_{F'}\Omega_{F'/\kappa}
       \longrightarrow\Omega_{S/\kappa}
       \longrightarrow\Omega_{S/F'}\longrightarrow0.
\]
The first two spaces both have dimension $n$ over $S$, since the
fields are finitely generated of transcendence degree $n$ over the
perfect field $\kappa$.  The dimension of the kernel of the first
map is therefore $\dim_S\Omega_{S/F'}$.
Let $S_{\mathrm{sep}}/F'$ be the maximal separable subextension and
write $[S:S_{\mathrm{sep}}]=p^u$.  By
\eqref{eq:compositum-small-p-part}, $u\leq h$.  The purely
inseparable extension $S/S_{\mathrm{sep}}$ can be generated by at
most $u$ elements: every nontrivial successive adjunction has degree
at least $p$.  Their differentials generate
$\Omega_{S/S_{\mathrm{sep}}}$, while
$\Omega_{S_{\mathrm{sep}}/F'}=0$.  The differential sequence thus
gives $\dim_S\Omega_{S/F'}\leq u\leq h$.

Any subspace spanned by $|I|$ of the original basis vectors loses
at most $h$ dimensions under a map with kernel dimension at most
$h$.  This proves the assertion for the forms.  Projection of
\eqref{eq:label-covectors} onto the second summand gives their
negatives, so the rank of the covectors is at least the rank of the
forms.
\end{proof}

The extra coordinate in \eqref{eq:label-covectors} records the
coefficient of the elliptic curve's invariant differential. Indeed,
on the constant curve $A_{0,S}$, absolute cotangents split into the
line generated by $\eta$ and the constant differentials from $S$.
Restriction along the marked section $t_i$ is therefore
\[
 S\oplus\Omega_{S/\kappa}\longrightarrow\Omega_{S/\kappa},
 \qquad (\lambda,\beta)\longmapsto\lambda t_i^*\eta+\beta.
\]
Its kernel is the line generated by $(1,-t_i^*\eta)$. These
cotangent directions normal to the marked sections will give the
ramification constraints in \Cref{sec:differentials}.
